\chapter{Вплив матеріальної неоднорідності за заданої геометрії тріщини}
\label{ch:fixed-crack}

Перш ніж будувати траєкторії тріщини в метакомпозиті, необхідно встановити, як зміна пружних характеристик кілець впливає на розв'язок при незмінній геометрії тріщини. Це дає змогу відокремити прямий матеріальний ефект від зміни розв'язку, спричиненої поворотом та подовженням тріщини. У цьому розділі розглядаємо найпростіший варіант матеріальної архітектури: модулі всіх трьох кілець змінюються одночасно і в однаковій пропорції. З'ясуємо, чи зберігаються знак і величина впливу кілець при наближенні вершини до області неоднорідності, а також чи можна передбачити цей вплив за зміною інтегральної жорсткості пластини.

\section{Схема розрахункового експерименту}
\label{sec:fixed-design}

Вихідними є три геометрії тріщини, отримані для однорідної перфорованої пластини: початкова пряма тріщина S0 та дві ламані після одного і двох приростів за критерієм МКН (S1 і S2). Довжини тріщин становлять 25,4; 29,4 та 33,4 мм. На кожній з цих геометрій розв'язуємо крайову задачу для трьох розподілів модулів:
\begin{equation}
 \boldsymbol\lambda_{\mathrm H}=(1,1,1),\qquad
 \boldsymbol\lambda_{\mathrm s}=(0{,}5,0{,}5,0{,}5),\qquad
 \boldsymbol\lambda_{\mathrm r}=(2,2,2).
 \label{eq:fixed-architectures}
\end{equation}
Позначення s і r у нижніх індексах стосуються відповідно м'яких і жорстких кілець; аргумент $s$ надалі означає довжину тріщини. Ці конфігурації з трьома активними кільцями не слід ототожнювати з soft1, soft2 та stiff13, які мають інше розташування активних областей.

У межах кожної геометрії використовуємо однакову сітку, крайові умови та силу $P=1$ Н. Змінюємо лише матеріальні характеристики елементів кілець. Однорідний розв'язок, отриманий за допомогою багатофазної реалізації при $\lambda_i=1$, побітово відтворює відповідний однорідний розв'язок. Отже, мітки матеріальних областей самі по собі не створюють чисельного контрасту. Усі три геометрії мають однакову дискретну площу кільця $1012{,}152529$ мм$^2$ та 76 хорд на його зовнішній межі. Однак зовнішня сітка змінюється від стану до стану, а S0 належить до початкового сімейства сіток прямої тріщини.

\begin{figure}[htbp]
 \centering
 % Nominal rings; reference tips from WP34_STATE2_RESULTS.md.
\begin{tikzpicture}[x=.7mm,y=.7mm,font=\small]
 \draw (65,0)--(205,0);
 \foreach \yy in {69.85,120.65}{
   \filldraw[fill=black!8,draw=black!55] (152.4,\yy) circle[radius=19.05];
   \filldraw[fill=white,draw=black] (152.4,\yy) circle[radius=6.35];
 }
 \node[anchor=west] at (174,69.85) {Кільце 1};
 \node[anchor=west] at (174,120.65) {Кільце 2};
 \draw[very thick] (101.6,0)--(101.6,25.4)--(102.4722795671,29.3037326185)--(103.4091509165,33.1924691839);
 \foreach \xx/\yy/\lab/\ly in {101.6/25.4/S0/16,102.4722795671/29.3037326185/S1/29,103.4091509165/33.1924691839/S2/42}{
   \fill (\xx,\yy) circle[radius=1.1pt];
   \draw[black!65] (\xx,\yy)--(89,\ly)--(81,\ly);
   \node[anchor=east] at (80,\ly) {\lab};
 }
 \draw[densely dashed] (101.6,25.4)--(138.06341397,57.30548722);
 \draw[densely dashed] (103.4091509165,33.1924691839)--(137.14721633,58.43706477);
 \node[anchor=west,align=left,font=\scriptsize] at (110,18) {Відстані до найближчого\\вузла кільця: 48,46;\\45,27; 42,14 мм};
 \draw[->] (66,2)--(79,2) node[anchor=west] {$x$};
 \draw[->] (66,2)--(66,15) node[anchor=south] {$y$};
 \node[anchor=north] at (135,-3) {Фрагмент пластини; розміри в мм};
\end{tikzpicture}

 \caption{Три зафіксовані геометрії тріщини в однорідній пластині та їхнє положення відносно нижніх кілець. Номінальні кругові контури й штрихові лінії показано схематично; відстані наведено за фактичними вузлами сітки.}
 \label{fig:fixed-geometry}
\end{figure}

Як видно з рис.~\ref{fig:fixed-geometry}, тріщина ще не наближається до міжфазної межі настільки, щоб вона впливала на допустимість привершинної області. Відстань до найближчого вузла кільця зменшується з 48,46 до 42,14 мм, а запас відносно номінального зовнішнього радіуса інтегрування 2,6 мм залишається більшим за 39 мм. Тому розглянуті зміни КІН є результатом перерозподілу пружних полів у пластині, а не зміни локального матеріалу вершини.

\section{Екранування та підсилення рушійної сили}
\label{sec:fixed-energy}

Відносний ефект визначаємо за~\eqref{eq:problem-fixed-contrast}. Для зручності позначимо $R_{\mathrm{mat}}=1+F_G^{\mathrm{fix}}$. У таблиці~\ref{tab:fixed-energy} наведено основні EDI-значення контрастів та два прийняті показники чисельної невизначеності, визначені в розділі~\ref{ch:verification}. Усі значення $F_G$ наведено у відсотках, а $U_{\mathrm{prim}}$ та $U_{\mathrm{cons}}$ --- у безрозмірних одиницях.

\begin{table}[htbp]
 \centering
 \caption{Прямий матеріальний ефект для трьох зафіксованих геометрій}
 \label{tab:fixed-energy}
 \begin{tabular}{llrrr}
 \toprule
 Стан & Кільця & $F_G$, \% & $U_{\mathrm{prim}}$ & $U_{\mathrm{cons}}$\\
 \midrule
 S0 & М'які & $-0,6093$ & $9,385\cdot10^{-6}$ & 0,0349\\
    & Жорсткі & $-0,9040$ & $6,839\cdot10^{-6}$ & 0,0349\\
 S1 & М'які & $+0,6199$ & $4,699\cdot10^{-6}$ & 0,0352\\
    & Жорсткі & $-1,8130$ & $2,586\cdot10^{-6}$ & 0,0344\\
 S2 & М'які & $+2,1813$ & $1,303\cdot10^{-6}$ & 0,0358\\
    & Жорсткі & $-2,9313$ & $4,968\cdot10^{-8}$ & 0,0340\\
 \bottomrule
 \end{tabular}
\end{table}

\begin{figure}[htbp]
 \centering
 \begin{tikzpicture}[x=1mm,y=1mm,font=\small,line width=.45pt]
\draw[black!15] (0,4.3077)--(115,4.3077);\node[anchor=east,font=\scriptsize] at (-1,4.3077) {$-6$};
\draw[black!15] (0,12.9231)--(115,12.9231);\node[anchor=east,font=\scriptsize] at (-1,12.9231) {$-4$};
\draw[black!15] (0,21.5385)--(115,21.5385);\node[anchor=east,font=\scriptsize] at (-1,21.5385) {$-2$};
\draw[black!15] (0,30.1538)--(115,30.1538);\node[anchor=east,font=\scriptsize] at (-1,30.1538) {$0$};
\draw[black!15] (0,38.7692)--(115,38.7692);\node[anchor=east,font=\scriptsize] at (-1,38.7692) {$2$};
\draw[black!15] (0,47.3846)--(115,47.3846);\node[anchor=east,font=\scriptsize] at (-1,47.3846) {$4$};
\draw[black!15] (0,56)--(115,56);\node[anchor=east,font=\scriptsize] at (-1,56) {$6$};
\draw[densely dotted] (0,30.1538)--(115,30.1538);
\draw (0,0) rectangle (115,56);\node[anchor=south] at (57.5,57) {$F_G$, \%};
\draw (8,0)--(8,-1);\node[anchor=north,font=\scriptsize] at (8,-1.5) {S0};
\draw (57.5,0)--(57.5,-1);\node[anchor=north,font=\scriptsize] at (57.5,-1.5) {S1};
\draw (107,0)--(107,-1);\node[anchor=north,font=\scriptsize] at (107,-1.5) {S2};
\draw[blue!65!black,opacity=.58] (6.2,12.4953)--(6.2,42.563) (4.9,12.4953)--(7.5,12.4953) (4.9,42.563)--(7.5,42.563);
\draw[blue!65!black,opacity=.58] (55.7,17.661)--(55.7,47.9872) (54.4,17.661)--(57,17.661) (54.4,47.9872)--(57,47.9872);
\draw[blue!65!black,opacity=.58] (105.2,24.1288)--(105.2,54.9718) (103.9,24.1288)--(106.5,24.1288) (103.9,54.9718)--(106.5,54.9718);
\draw[blue!65!black,line width=.8pt] (6.2,27.5291)--(55.7,32.8241)--(105.2,39.5503);
\draw[blue!65!black,fill=blue!65!black] (5.5,26.8291) rectangle (6.9,28.2291);
\draw[blue!65!black,fill=blue!65!black] (55,32.1241) rectangle (56.4,33.5241);
\draw[blue!65!black,fill=blue!65!black] (104.5,38.8503) rectangle (105.9,40.2503);
\draw[blue!65!black,fill=white] (108.2,40.6566)--(109.3,39.5566)--(108.2,38.4566)--(107.1,39.5566)--cycle;
\draw[red!65!black,opacity=.58] (9.8,11.2258)--(9.8,41.2935) (8.5,11.2258)--(11.1,11.2258) (8.5,41.2935)--(11.1,41.2935);
\draw[red!65!black,opacity=.58] (59.3,7.5256)--(59.3,37.1625) (58,7.5256)--(60.6,7.5256) (58,37.1625)--(60.6,37.1625);
\draw[red!65!black,opacity=.58] (108.8,2.8807)--(108.8,32.173) (107.5,2.8807)--(110.1,2.8807) (107.5,32.173)--(110.1,32.173);
\draw[red!65!black,line width=.8pt] (9.8,26.2597)--(59.3,22.3441)--(108.8,17.5268);
\draw[red!65!black,fill=red!65!black] (9.8,27.2097)--(8.85,25.6097)--(10.75,25.6097)--cycle;
\draw[red!65!black,fill=red!65!black] (59.3,23.2941)--(58.35,21.6941)--(60.25,21.6941)--cycle;
\draw[red!65!black,fill=red!65!black] (108.8,18.4768)--(107.85,16.8768)--(109.75,16.8768)--cycle;
\draw[red!65!black,fill=white] (111.8,18.6183)--(112.9,17.5183)--(111.8,16.4183)--(110.7,17.5183)--cycle;
\draw[blue!65!black,line width=.8pt] (5,-12)--(14,-12);
\draw[blue!65!black,fill=blue!65!black] (8.8,-12.7) rectangle (10.2,-11.3);
\node[anchor=west,font=\scriptsize] at (15,-12) {М'які кільця};
\draw[red!65!black,line width=.8pt] (51,-12)--(60,-12);
\draw[red!65!black,fill=red!65!black] (55.5,-11.05)--(54.55,-12.65)--(56.45,-12.65)--cycle;
\node[anchor=west,font=\scriptsize] at (61,-12) {Жорсткі кільця};
\node[anchor=west,font=\scriptsize] at (0,-19) {Ромби: S2alt; відрізки: $\pm U_{\mathrm{cons}}$.};
\end{tikzpicture}

 \caption{Енергетичний контраст за незмінної геометрії тріщини. Відрізки показують консервативний показник $\pm U_{\mathrm{cons}}$, а не довірчі інтервали; первинний розкид наведено в табл.~\ref{tab:fixed-energy}. Лінії лише з'єднують дискретні результати; ромби відповідають альтернативній сітці S2alt.}
 \label{fig:fixed-contrast}
\end{figure}

При збільшенні модуля Юнга кілець у два рази отримано від'ємні контрасти на всіх трьох геометріях, причому їхні модулі зростають. При зменшенні модуля вдвічі отримано іншу поведінку: від'ємне значення на S0 змінюється додатним на S1 і зростає на S2. Таким чином, для однієї і тієї самої неоднорідності знак обчисленої зміни рушійної сили залежить від геометрії тріщини. За трьома дискретними точками не можна визначити точну довжину, за якої відбувається зміна знака, або довести монотонність між ними.

Усі шість контрастів розрізняються за первинним критерієм $|F_G|>2U_{\mathrm{prim}}$, причому знаки за COD та EDI збігаються. Проте жоден не розрізняється за консервативним критерієм $|F_G|>U_{\mathrm{cons}}$. Це стосується саме раннього дослідження трьох однаково активованих кілець. На S2 додатково перевірено обидва контрасти на незалежно перебудованій зовнішній сітці з 112 хордами: для м'яких кілець $+2{,}1813\to+2{,}1828$\,\%, для жорстких $-2{,}9313\to-2{,}9333$\,\%. Спостережена зміна становить близько 0,07\,\% від величини відповідного ефекту. Цей результат підтримує стійкість контрастів на S2, але не є прямою перевіркою сіткової чутливості для S0 і S1 та не скасовує консервативного висновку.

\section{Внески двох мод руйнування}
\label{sec:fixed-modes}

Щоб зрозуміти причину зміни енергетичного контрасту, розглянемо окремі внески КІН. Для будь-якої з фіксованих геометрій позначимо однорідні КІН через $A=K_{\mathrm I,H}$ та $B_K=K_{\mathrm{II},H}$, а зміни в неоднорідному матеріалі --- через $\Delta\KI$ та $\Delta\KII$. Оскільки матеріал вершини незмінний, з~\eqref{eq:problem-G} випливає точна тотожність
\begin{equation}
 \begin{aligned}
 F_G={}&F_{\mathrm I}+F_{\mathrm{II}},\\
 F_{\mathrm I}={}&\frac{2A\Delta\KI+(\Delta\KI)^2}{A^2+B_K^2},\\
 F_{\mathrm{II}}={}&\frac{2B_K\Delta\KII+(\Delta\KII)^2}{A^2+B_K^2}.
 \end{aligned}
 \label{eq:fixed-mode-decomp}
\end{equation}
Тут $F_{\mathrm I}$ та $F_{\mathrm{II}}$ є змінами відповідних модових складових енергії, нормованими на повне опорне $G$. Це алгебраїчний розклад, а не розділення просторових механізмів перерозподілу напружень.

\begin{figure}[htbp]
 \centering
 \begin{tikzpicture}[x=1mm,y=1mm,font=\small,line width=.45pt]
\draw[black!15] (0,0)--(65,0);\node[anchor=east,font=\scriptsize] at (-1,0) {$0{,}15$};
\draw[black!15] (0,13.6667)--(65,13.6667);\node[anchor=east,font=\scriptsize] at (-1,13.6667) {$0{,}17$};
\draw[black!15] (0,27.3333)--(65,27.3333);\node[anchor=east,font=\scriptsize] at (-1,27.3333) {$0{,}19$};
\draw[black!15] (0,41)--(65,41);\node[anchor=east,font=\scriptsize] at (-1,41) {$0{,}21$};
\draw (0,0) rectangle (65,41);\node[anchor=south] at (32.5,42) {(а) $K_{\mathrm I}$, МПа$\sqrt{\text{м}}$/кН};
\draw (8,0)--(8,-1);\node[anchor=north,font=\scriptsize] at (8,-1.5) {S0};
\draw (32.5,0)--(32.5,-1);\node[anchor=north,font=\scriptsize] at (32.5,-1.5) {S1};
\draw (57,0)--(57,-1);\node[anchor=north,font=\scriptsize] at (57,-1.5) {S2};
\draw[black,line width=.8pt] (8,11.2312)--(32.5,22.2663)--(57,31.4213);
\draw[black,fill=black] (8,11.2312) circle[radius=.7mm];
\draw[black,fill=black] (32.5,22.2663) circle[radius=.7mm];
\draw[black,fill=black] (57,31.4213) circle[radius=.7mm];
\draw[blue!65!black,line width=.8pt] (8,10.9597)--(32.5,22.6565)--(57,32.8847);
\draw[blue!65!black,fill=blue!65!black] (7.3,10.2597) rectangle (8.7,11.6597);
\draw[blue!65!black,fill=blue!65!black] (31.8,21.9565) rectangle (33.2,23.3565);
\draw[blue!65!black,fill=blue!65!black] (56.3,32.1847) rectangle (57.7,33.5847);
\draw[red!65!black,line width=.8pt] (8,10.6639)--(32.5,21.1246)--(57,29.4345);
\draw[red!65!black,fill=red!65!black] (8,11.6139)--(7.05,10.0139)--(8.95,10.0139)--cycle;
\draw[red!65!black,fill=red!65!black] (32.5,22.0746)--(31.55,20.4746)--(33.45,20.4746)--cycle;
\draw[red!65!black,fill=red!65!black] (57,30.3845)--(56.05,28.7845)--(57.95,28.7845)--cycle;
\draw[black,densely dashed] (8,10.234)--(32.5,21.1692)--(57,30.2443);
\draw[black,fill=white] (8,10.234) circle[radius=.7mm];
\draw[black,fill=white] (32.5,21.1692) circle[radius=.7mm];
\draw[black,fill=white] (57,30.2443) circle[radius=.7mm];
\draw[blue!65!black,densely dashed] (8,9.9648)--(32.5,21.5556)--(57,31.6947);
\draw[blue!65!black,fill=white] (7.3,9.2648) rectangle (8.7,10.6648);
\draw[blue!65!black,fill=white] (31.8,20.8556) rectangle (33.2,22.2556);
\draw[blue!65!black,fill=white] (56.3,30.9947) rectangle (57.7,32.3947);
\draw[red!65!black,densely dashed] (8,9.6718)--(32.5,20.0377)--(57,28.275);
\draw[red!65!black,fill=white] (8,10.6218)--(7.05,9.0218)--(8.95,9.0218)--cycle;
\draw[red!65!black,fill=white] (32.5,20.9877)--(31.55,19.3877)--(33.45,19.3877)--cycle;
\draw[red!65!black,fill=white] (57,29.225)--(56.05,27.625)--(57.95,27.625)--cycle;
\draw[black!15] (82,0)--(147,0);\node[anchor=east,font=\scriptsize] at (81,0) {$0$};
\draw[black!15] (82,9.7619)--(147,9.7619);\node[anchor=east,font=\scriptsize] at (81,9.7619) {$5$};
\draw[black!15] (82,19.5238)--(147,19.5238);\node[anchor=east,font=\scriptsize] at (81,19.5238) {$10$};
\draw[black!15] (82,29.2857)--(147,29.2857);\node[anchor=east,font=\scriptsize] at (81,29.2857) {$15$};
\draw[black!15] (82,39.0476)--(147,39.0476);\node[anchor=east,font=\scriptsize] at (81,39.0476) {$20$};
\draw (82,0) rectangle (147,41);\node[anchor=south] at (114.5,42) {(б) $10^3 K_{\mathrm{II}}$, МПа$\sqrt{\text{м}}$/кН};
\draw (90,0)--(90,-1);\node[anchor=north,font=\scriptsize] at (90,-1.5) {S0};
\draw (114.5,0)--(114.5,-1);\node[anchor=north,font=\scriptsize] at (114.5,-1.5) {S1};
\draw (139,0)--(139,-1);\node[anchor=north,font=\scriptsize] at (139,-1.5) {S2};
\draw[black,line width=.8pt] (90,36.7576)--(114.5,2.9552)--(139,6.1499);
\draw[black,fill=black] (90,36.7576) circle[radius=.7mm];
\draw[black,fill=black] (114.5,2.9552) circle[radius=.7mm];
\draw[black,fill=black] (139,6.1499) circle[radius=.7mm];
\draw[blue!65!black,line width=.8pt] (90,34.6852)--(114.5,0.6467)--(139,3.8706);
\draw[blue!65!black,fill=blue!65!black] (89.3,33.9852) rectangle (90.7,35.3852);
\draw[blue!65!black,fill=blue!65!black] (113.8,-0.0533) rectangle (115.2,1.3467);
\draw[blue!65!black,fill=blue!65!black] (138.3,3.1706) rectangle (139.7,4.5706);
\draw[red!65!black,line width=.8pt] (90,37.883)--(114.5,4.4478)--(139,7.5477);
\draw[red!65!black,fill=red!65!black] (90,38.833)--(89.05,37.233)--(90.95,37.233)--cycle;
\draw[red!65!black,fill=red!65!black] (114.5,5.3978)--(113.55,3.7978)--(115.45,3.7978)--cycle;
\draw[red!65!black,fill=red!65!black] (139,8.4977)--(138.05,6.8977)--(139.95,6.8977)--cycle;
\draw[black,densely dashed] (90,36.2526)--(114.5,2.8954)--(139,6.0556);
\draw[black,fill=white] (90,36.2526) circle[radius=.7mm];
\draw[black,fill=white] (114.5,2.8954) circle[radius=.7mm];
\draw[black,fill=white] (139,6.0556) circle[radius=.7mm];
\draw[blue!65!black,densely dashed] (90,34.207)--(114.5,0.6162)--(139,3.8053);
\draw[blue!65!black,fill=white] (89.3,33.507) rectangle (90.7,34.907);
\draw[blue!65!black,fill=white] (113.8,-0.0838) rectangle (115.2,1.3162);
\draw[blue!65!black,fill=white] (138.3,3.1053) rectangle (139.7,4.5053);
\draw[red!65!black,densely dashed] (90,37.3637)--(114.5,4.3693)--(139,7.436);
\draw[red!65!black,fill=white] (90,38.3137)--(89.05,36.7137)--(90.95,36.7137)--cycle;
\draw[red!65!black,fill=white] (114.5,5.3193)--(113.55,3.7193)--(115.45,3.7193)--cycle;
\draw[red!65!black,fill=white] (139,8.386)--(138.05,6.786)--(139.95,6.786)--cycle;
\draw[black,line width=.8pt] (5,-12)--(14,-12);
\draw[black,fill=black] (9.5,-12) circle[radius=.7mm];
\node[anchor=west,font=\scriptsize] at (15,-12) {H};
\draw[blue!65!black,line width=.8pt] (51,-12)--(60,-12);
\draw[blue!65!black,fill=blue!65!black] (54.8,-12.7) rectangle (56.2,-11.3);
\node[anchor=west,font=\scriptsize] at (61,-12) {М'які кільця};
\draw[red!65!black,line width=.8pt] (97,-12)--(106,-12);
\draw[red!65!black,fill=red!65!black] (101.5,-11.05)--(100.55,-12.65)--(102.45,-12.65)--cycle;
\node[anchor=west,font=\scriptsize] at (107,-12) {Жорсткі кільця};
\node[anchor=west,font=\scriptsize] at (0,-19) {Суцільні лінії та заповнені маркери: EDI; штрихові та порожні: COD.};
\end{tikzpicture}

 \caption{Знакові КІН двох мод для трьох геометрій. EDI та COD показано окремо. Значення нормовано на силу 1 кН; на панелі (б) числові значення помножено на $10^3$.}
 \label{fig:fixed-K}
\end{figure}

\begin{table}[htbp]
 \centering
 \caption{Внески мод I та II до зміни рушійної сили (у відсотках від опорного $G$)}
 \label{tab:fixed-modes}
 \begin{tabular}{llrrr}
 \toprule
 Стан & Кільця & $F_{\mathrm I}$ & $F_{\mathrm{II}}$ & $F_G$\\
 \midrule
S0 & М'які & $-0,4709$ & $-0,1385$ & $-0,6093$ \\
S0 & Жорсткі & $-0,9825$ & $0,0785$ & $-0,9040$ \\
S1 & М'які & $0,6264$ & $-0,0065$ & $0,6199$ \\
S1 & Жорсткі & $-1,8217$ & $0,0087$ & $-1,8130$ \\
S2 & М'які & $2,1969$ & $-0,0156$ & $2,1813$ \\
S2 & Жорсткі & $-2,9443$ & $0,0131$ & $-2,9313$ \\
 \bottomrule
 \end{tabular}
\end{table}

Для м'яких кілець на S0 обидві складові від'ємні; зменшення КІН моди II забезпечує близько 23\,\% загального від'ємного контрасту. На S1 і S2 внесок моди I змінює знак, тоді як внесок моди II залишається від'ємним. Отже, обчислений перехід від екранування до підсилення визначається саме зміною КІН моди I, а не зміною знака $\KII$. Для жорстких кілець, навпаки, внесок моди I від'ємний, а моди II додатний на всіх трьох геометріях. Після першого повороту тріщини енергетична зміна майже повністю визначається модою I, хоча модові коефіцієнти II залишаються важливими для визначення напрямку.

Величина $\KII$ на S1 в однорідній пластині становить лише $0{,}001514$ МПа$\sqrt{\text{м}}$ на 1 кН порівняно з $0{,}018827$ на S0. Тому відносна зміна $\KII$ при заміні матеріалу може бути великою навіть за близької абсолютної зміни. Для фізичного тлумачення напрямних ефектів доцільніше розглядати знакові абсолютні зміни та відношення $\KII/\KI$, а не лише відносну зміну другого КІН.

\section{Вплив кілець на локальний напрямок приросту}
\label{sec:fixed-steering}

Позначимо $p=\KII/\KI$ та $\psi=\arctan p$. На кожній геометрії обчислюємо кут МКН за~\eqref{eq:method-mts} в одному й тому самому локальному базисі для всіх трьох матеріалів. Різниця кутів
\begin{equation}
 \Delta\theta=\theta_X-\theta_{\mathrm H}
 \label{eq:fixed-dtheta}
\end{equation}
є прогнозом зміни напрямку наступного приросту, а не різницею вже сформованих траєкторій.

\begin{figure}[htbp]
 \centering
 \begin{tikzpicture}[x=1mm,y=1mm,font=\small,line width=.45pt]
\draw[black!15] (0,0)--(65,0);\node[anchor=east,font=\scriptsize] at (-1,0) {$-0{,}6$};
\draw[black!15] (0,8.2)--(65,8.2);\node[anchor=east,font=\scriptsize] at (-1,8.2) {$-0{,}3$};
\draw[black!15] (0,16.4)--(65,16.4);\node[anchor=east,font=\scriptsize] at (-1,16.4) {$0$};
\draw[black!15] (0,24.6)--(65,24.6);\node[anchor=east,font=\scriptsize] at (-1,24.6) {$0{,}3$};
\draw[black!15] (0,32.8)--(65,32.8);\node[anchor=east,font=\scriptsize] at (-1,32.8) {$0{,}6$};
\draw[black!15] (0,41)--(65,41);\node[anchor=east,font=\scriptsize] at (-1,41) {$0{,}9$};
\draw[densely dotted] (0,16.4)--(65,16.4);
\draw (0,0) rectangle (65,41);\node[anchor=south] at (32.5,42) {(а) $\Delta\theta$, град};
\draw (8,0)--(8,-1);\node[anchor=north,font=\scriptsize] at (8,-1.5) {S0};
\draw (32.5,0)--(32.5,-1);\node[anchor=north,font=\scriptsize] at (32.5,-1.5) {S1};
\draw (57,0)--(57,-1);\node[anchor=north,font=\scriptsize] at (57,-1.5) {S2};
\draw[blue!65!black,line width=.8pt] (8,34.0819)--(32.5,36.6975)--(57,35.3779);
\draw[blue!65!black,fill=blue!65!black] (7.3,33.3819) rectangle (8.7,34.7819);
\draw[blue!65!black,fill=blue!65!black] (31.8,35.9975) rectangle (33.2,37.3975);
\draw[blue!65!black,fill=blue!65!black] (56.3,34.6779) rectangle (57.7,36.0779);
\draw[red!65!black,line width=.8pt] (8,4.7956)--(32.5,2.9354)--(57,4.056);
\draw[red!65!black,fill=red!65!black] (8,5.7456)--(7.05,4.1456)--(8.95,4.1456)--cycle;
\draw[red!65!black,fill=red!65!black] (32.5,3.8854)--(31.55,2.2854)--(33.45,2.2854)--cycle;
\draw[red!65!black,fill=red!65!black] (57,5.006)--(56.05,3.406)--(57.95,3.406)--cycle;
\draw[black!15] (82,0)--(147,0);\node[anchor=east,font=\scriptsize] at (81,0) {$0$};
\draw[black!15] (82,11.7143)--(147,11.7143);\node[anchor=east,font=\scriptsize] at (81,11.7143) {$2$};
\draw[black!15] (82,23.4286)--(147,23.4286);\node[anchor=east,font=\scriptsize] at (81,23.4286) {$4$};
\draw[black!15] (82,35.1429)--(147,35.1429);\node[anchor=east,font=\scriptsize] at (81,35.1429) {$6$};
\draw (82,0) rectangle (147,41);\node[anchor=south] at (114.5,42) {(б) $\psi$, град};
\draw (90,0)--(90,-1);\node[anchor=north,font=\scriptsize] at (90,-1.5) {S0};
\draw (114.5,0)--(114.5,-1);\node[anchor=north,font=\scriptsize] at (114.5,-1.5) {S1};
\draw (139,0)--(139,-1);\node[anchor=north,font=\scriptsize] at (139,-1.5) {S2};
\draw[black,line width=.8pt] (90,37.8009)--(114.5,2.782)--(139,5.3934);
\draw[black,fill=black] (90,37.8009) circle[radius=.7mm];
\draw[black,fill=black] (114.5,2.782) circle[radius=.7mm];
\draw[black,fill=black] (139,5.3934) circle[radius=.7mm];
\draw[blue!65!black,line width=.8pt] (90,35.7709)--(114.5,0.6069)--(139,3.3579);
\draw[blue!65!black,fill=blue!65!black] (89.3,35.0709) rectangle (90.7,36.4709);
\draw[blue!65!black,fill=blue!65!black] (113.8,-0.0931) rectangle (115.2,1.3069);
\draw[blue!65!black,fill=blue!65!black] (138.3,2.6579) rectangle (139.7,4.0579);
\draw[red!65!black,line width=.8pt] (90,39.1415)--(114.5,4.2256)--(139,6.7185);
\draw[red!65!black,fill=red!65!black] (90,40.0915)--(89.05,38.4915)--(90.95,38.4915)--cycle;
\draw[red!65!black,fill=red!65!black] (114.5,5.1756)--(113.55,3.5756)--(115.45,3.5756)--cycle;
\draw[red!65!black,fill=red!65!black] (139,7.6685)--(138.05,6.0685)--(139.95,6.0685)--cycle;
\draw[black,line width=.8pt] (5,-12)--(14,-12);
\draw[black,fill=black] (9.5,-12) circle[radius=.7mm];
\node[anchor=west,font=\scriptsize] at (15,-12) {H};
\draw[blue!65!black,line width=.8pt] (51,-12)--(60,-12);
\draw[blue!65!black,fill=blue!65!black] (54.8,-12.7) rectangle (56.2,-11.3);
\node[anchor=west,font=\scriptsize] at (61,-12) {М'які кільця};
\draw[red!65!black,line width=.8pt] (97,-12)--(106,-12);
\draw[red!65!black,fill=red!65!black] (101.5,-11.05)--(100.55,-12.65)--(102.45,-12.65)--cycle;
\node[anchor=west,font=\scriptsize] at (107,-12) {Жорсткі кільця};
\end{tikzpicture}

 \caption{Різниця прогнозованих кутів МКН відносно однорідної пластини (а) та змішаність мод $\psi$ (б) для фіксованих геометрій.}
 \label{fig:fixed-direction}
\end{figure}

На всіх трьох геометріях м'які кільця зменшують модуль від'ємного кута МКН, а жорсткі його збільшують. Для м'яких кілець $\Delta\theta=+0{,}647^\circ$, $+0{,}743^\circ$ та $+0{,}694^\circ$; для жорстких відповідно $-0{,}425^\circ$, $-0{,}493^\circ$ та $-0{,}452^\circ$. Усі шість різниць розрізняються за прийнятим первинним кутовим критерієм, для якого розкид не перевищує $0{,}003^\circ$. На S2 зміна $\Delta\theta$ при заміні зовнішньої сітки не перевищує $2\cdot10^{-4}$ градуса. Це спостереження стосується прогнозу на заданій тріщині.

Математична причина відмінності енергетичного й напрямного ефектів випливає з різної залежності $G$ і $\theta$ від двох КІН. Диференціюючи енергетичне співвідношення при незмінному модулі матриці, отримуємо
\begin{equation}
 \frac{\mathrm dG}{G}=
 \frac{2(\KI\,\mathrm d\KI+\KII\,\mathrm d\KII)}
 {\KI^2+\KII^2}.
 \label{eq:fixed-G-differential}
\end{equation}
Умова стаціонарності МКН~\eqref{eq:method-mts-stationarity} після ділення на $\KI>0$ дає
\begin{equation}
 \frac{\mathrm d\theta}{\mathrm dp}
 =-\frac{3\cos\theta-1}
 {\cos\theta-3p\sin\theta},\qquad
 \mathrm dp=\frac{\mathrm d\KII-p\,\mathrm d\KI}{\KI}.
 \label{eq:fixed-theta-derivative}
\end{equation}
Ці вирази справедливі на регулярній гілці максимуму, де знаменник ненульовий; кут вимірюється в радіанах. Для $|p|\ll1$
\begin{equation}
 \mathrm d\theta\simeq
 -\frac{2}{\KI}(\mathrm d\KII-p\,\mathrm d\KI),\qquad
 \frac{\mathrm dG}{G}\simeq
 \frac{2}{\KI}(\mathrm d\KI+p\,\mathrm d\KII).
 \label{eq:fixed-small-mixity}
\end{equation}
Звідси видно, що за малої змішаності енергетична характеристика переважно чутлива до зміни $\KI$, а напрямна --- до зміни $\KII$. Це локальна математична властивість двох функціоналів, а не твердження про незалежне керування ними через модулі кілець.

Для м'яких кілець абсолютна зміна $\KII$ від S0 до S2 становить $-1{,}0615$, $-1{,}1824$ та $-1{,}1675$ в одиницях $10^{-3}$ МПа$\sqrt{\text{м}}$/кН. Для жорстких --- $+0{,}5764$, $+0{,}7645$ та $+0{,}7159$. Ці зміни зберігають знак, тоді як зміна $\KI$ для м'яких кілець змінює знак. Отже, спостережувана зміна знака енергетичного ефекту при незмінному знаку напрямного ефекту узгоджується з математичною структурою~\eqref{eq:fixed-G-differential} та~\eqref{eq:fixed-theta-derivative}. Причину самого перерозподілу КІН необхідно шукати в полі напружень неоднорідної пластини.

\section{Податливість пластини та енергетична рушійна сила}
\label{sec:fixed-compliance}

Введемо податливість $C=-u_y(\boldsymbol x_P)/P$ за додатної величини спрямованої вниз сили. За незмінної геометрії й лінійно-пружної поведінки повна енергія деформації дорівнює $\mathcal U=P^2C/2$. Якщо $s$ --- довжина тріщини, а похідну взято за її віртуального прямолінійного приросту в локальному напрямку, то при керуванні силою
\begin{equation}
 G=-\frac{1}{t}\frac{\partial\Pi}{\partial s}
 =\frac{P^2}{2t}\frac{\partial C}{\partial s}.
 \label{eq:fixed-compliance-G}
\end{equation}
Тут $\Pi=-P^2C/2$ --- повна потенціальна енергія рівноважної системи. Похідна має зміст для неперервної варіації геометрії при незмінному матеріальному розподілі; для точкового навантаження двовимірної моделі та дискретних розрахунків податливість у вузлі є сітково залежною. Співвідношення не використовується тут для незалежного точного обчислення $G$ різницями податливостей на перебудованих сітках.

Для фіксованої сітки можна встановити знак зміни $C$ без розв'язання задачі для кожного контрасту. Нехай $\boldsymbol K(\lambda)$ --- редукована матриця жорсткості з урахуванням товщини, $\boldsymbol f=P\boldsymbol e$, де $\boldsymbol e$ --- одиничний вектор навантаження, а $\lambda$ --- спільний множник модуля всіх кілець. За сталого коефіцієнта Пуассона
\begin{equation}
 \boldsymbol K(\lambda)=\boldsymbol K_{\mathrm m}
 +\lambda\boldsymbol K_{\mathrm r},\qquad
 C=\boldsymbol e^{\mathsf T}\boldsymbol K^{-1}\boldsymbol e.
 \label{eq:fixed-K-compliance}
\end{equation}
Тут $\boldsymbol K_{\mathrm r}$ --- внесок кілець при модулі матриці, додатно напіввизначена матриця. Диференціювання дає
\begin{equation}
 \frac{\partial C}{\partial\lambda}
 =-\boldsymbol e^{\mathsf T}\boldsymbol K^{-1}
 \boldsymbol K_{\mathrm r}\boldsymbol K^{-1}\boldsymbol e
 \le 0.
 \label{eq:fixed-C-monotonicity}
\end{equation}
Тому збільшення жорсткості кілець не може збільшити податливість при фіксованій геометрії та силовому навантаженні. Проте знак $\partial C/\partial\lambda$ не встановлює знака змішаної похідної $\partial^2 C/(\partial s\,\partial\lambda)$, яка визначає зміну $G$. Саме тут виникає принципова відмінність між глобальною жорсткістю та локальною характеристикою руйнування.

\begin{figure}[htbp]
 \centering
 \begin{tikzpicture}[x=1mm,y=1mm,font=\small,line width=.45pt]
\draw[black!15] (0,0)--(65,0);\node[anchor=east,font=\scriptsize] at (-1,0) {$8$};
\draw[black!15] (0,9.1111)--(65,9.1111);\node[anchor=east,font=\scriptsize] at (-1,9.1111) {$9$};
\draw[black!15] (0,18.2222)--(65,18.2222);\node[anchor=east,font=\scriptsize] at (-1,18.2222) {$10$};
\draw[black!15] (0,27.3333)--(65,27.3333);\node[anchor=east,font=\scriptsize] at (-1,27.3333) {$11$};
\draw[black!15] (0,36.4444)--(65,36.4444);\node[anchor=east,font=\scriptsize] at (-1,36.4444) {$12$};
\draw (0,0) rectangle (65,41);\node[anchor=south] at (32.5,42) {(а) $10^9G$, Н/мм};
\draw (8,0)--(8,-1);\node[anchor=north,font=\scriptsize] at (8,-1.5) {S0};
\draw (32.5,0)--(32.5,-1);\node[anchor=north,font=\scriptsize] at (32.5,-1.5) {S1};
\draw (57,0)--(57,-1);\node[anchor=north,font=\scriptsize] at (57,-1.5) {S2};
\draw[black,line width=.8pt] (8,5.3261)--(32.5,20.0577)--(57,34.2187);
\draw[black,fill=black] (8,5.3261) circle[radius=.7mm];
\draw[black,fill=black] (32.5,20.0577) circle[radius=.7mm];
\draw[black,fill=black] (57,34.2187) circle[radius=.7mm];
\draw[blue!65!black,line width=.8pt] (8,4.8495)--(32.5,20.6338)--(57,36.5551);
\draw[blue!65!black,fill=blue!65!black] (7.3,4.1495) rectangle (8.7,5.5495);
\draw[blue!65!black,fill=blue!65!black] (31.8,19.9338) rectangle (33.2,21.3338);
\draw[blue!65!black,fill=blue!65!black] (56.3,35.8551) rectangle (57.7,37.2551);
\draw[red!65!black,line width=.8pt] (8,4.619)--(32.5,18.3726)--(57,31.0791);
\draw[red!65!black,fill=red!65!black] (8,5.569)--(7.05,3.969)--(8.95,3.969)--cycle;
\draw[red!65!black,fill=red!65!black] (32.5,19.3226)--(31.55,17.7226)--(33.45,17.7226)--cycle;
\draw[red!65!black,fill=red!65!black] (57,32.0291)--(56.05,30.4291)--(57.95,30.4291)--cycle;
\draw[black!15] (82,0)--(147,0);\node[anchor=east,font=\scriptsize] at (81,0) {$-2$};
\draw[black!15] (82,10.25)--(147,10.25);\node[anchor=east,font=\scriptsize] at (81,10.25) {$-1$};
\draw[black!15] (82,20.5)--(147,20.5);\node[anchor=east,font=\scriptsize] at (81,20.5) {$0$};
\draw[black!15] (82,30.75)--(147,30.75);\node[anchor=east,font=\scriptsize] at (81,30.75) {$1$};
\draw[black!15] (82,41)--(147,41);\node[anchor=east,font=\scriptsize] at (81,41) {$2$};
\draw[densely dotted] (82,20.5)--(147,20.5);
\draw (82,0) rectangle (147,41);\node[anchor=south] at (114.5,42) {(б) Зміна податливості та CMOD, \%};
\draw (90,0)--(90,-1);\node[anchor=north,font=\scriptsize] at (90,-1.5) {S0};
\draw (114.5,0)--(114.5,-1);\node[anchor=north,font=\scriptsize] at (114.5,-1.5) {S1};
\draw (139,0)--(139,-1);\node[anchor=north,font=\scriptsize] at (139,-1.5) {S2};
\draw[blue!65!black,line width=.8pt] (90,35.4344)--(114.5,35.3654)--(139,35.361);
\draw[blue!65!black,fill=blue!65!black] (89.3,34.7344) rectangle (90.7,36.1344);
\draw[blue!65!black,fill=blue!65!black] (113.8,34.6654) rectangle (115.2,36.0654);
\draw[blue!65!black,fill=blue!65!black] (138.3,34.661) rectangle (139.7,36.061);
\draw[red!65!black,line width=.8pt] (90,9.7081)--(114.5,9.6969)--(139,9.6302);
\draw[red!65!black,fill=red!65!black] (90,10.6581)--(89.05,9.0581)--(90.95,9.0581)--cycle;
\draw[red!65!black,fill=red!65!black] (114.5,10.6469)--(113.55,9.0469)--(115.45,9.0469)--cycle;
\draw[red!65!black,fill=red!65!black] (139,10.5802)--(138.05,8.9802)--(139.95,8.9802)--cycle;
\draw[blue!65!black,densely dashed] (90,5.0978)--(114.5,11.1473)--(139,17.6221);
\draw[blue!65!black,fill=white] (89.3,4.3978) rectangle (90.7,5.7978);
\draw[blue!65!black,fill=white] (113.8,10.4473) rectangle (115.2,11.8473);
\draw[blue!65!black,fill=white] (138.3,16.9221) rectangle (139.7,18.3221);
\draw[red!65!black,densely dashed] (90,24.0937)--(114.5,19.6071)--(139,14.8347);
\draw[red!65!black,fill=white] (90,25.0437)--(89.05,23.4437)--(90.95,23.4437)--cycle;
\draw[red!65!black,fill=white] (114.5,20.5571)--(113.55,18.9571)--(115.45,18.9571)--cycle;
\draw[red!65!black,fill=white] (139,15.7847)--(138.05,14.1847)--(139.95,14.1847)--cycle;
\draw[black,line width=.8pt] (5,-12)--(14,-12);
\draw[black,fill=black] (9.5,-12) circle[radius=.7mm];
\node[anchor=west,font=\scriptsize] at (15,-12) {H};
\draw[blue!65!black,line width=.8pt] (51,-12)--(60,-12);
\draw[blue!65!black,fill=blue!65!black] (54.8,-12.7) rectangle (56.2,-11.3);
\node[anchor=west,font=\scriptsize] at (61,-12) {М'які кільця};
\draw[red!65!black,line width=.8pt] (97,-12)--(106,-12);
\draw[red!65!black,fill=red!65!black] (101.5,-11.05)--(100.55,-12.65)--(102.45,-12.65)--cycle;
\node[anchor=west,font=\scriptsize] at (107,-12) {Жорсткі кільця};
\node[anchor=west,font=\scriptsize] at (0,-19) {На (б): суцільні лінії --- податливість, штрихові --- розкриття гирла.};
\end{tikzpicture}

 \caption{Абсолютне значення $G$ за сили 1 Н (а); відносні зміни податливості та нормального розкриття гирла тріщини (CMOD) відносно однорідної пластини (б).}
 \label{fig:fixed-structure}
\end{figure}

Розрахунки підтверджують передбачуваний знак зміни податливості: для м'яких кілець вона зростає приблизно на 1,45\,\%, для жорстких зменшується приблизно на 1,05\,\%. Ці контрасти майже однакові для трьох геометрій, хоча відповідні енергетичні контрасти істотно змінюються. Нормальне розкриття гирла також не відтворює закономірність зміни $G$: на S2 воно дещо зменшується для обох типів кілець. Отже, жодна з цих глобальних характеристик не є заміною КІН або $G$.

Окремо слід враховувати зміну опорного $G$ при подовженні тріщини. Для однорідної пластини $G$ зростає від $8{,}58457\cdot10^{-9}$ до $1{,}17557\cdot10^{-8}$ Н/мм (приблизно на 36,9\,\%). Абсолютне $G$ зростає також для обох неоднорідних конфігурацій. Від'ємний контраст означає меншу рушійну силу відносно однорідної пластини при тій самій довжині та геометрії, але не зменшення абсолютного $G$ зі збільшенням довжини. У цій задачі екранування не можна ототожнювати зі стабілізацією або зупинкою тріщини.

\section{Підсумок та перехід до вибіркового розташування кілець}
\label{sec:fixed-summary}

Перший розрахунковий експеримент дозволив встановити три закономірності в межах досліджених фіксованих геометрій. По-перше, знак прямого енергетичного ефекту м'яких кілець змінюється, тоді як знак напрямного ефекту зберігається. По-друге, після першого зламу зміна $G$ майже повністю визначається модою I, проте зміна напрямку залишається чутливою до моди II. По-третє, майже незмінний контраст податливості не визначає контрасту рушійної сили. Усі енергетичні ефекти ранньої серії розрізняються за первинним, але не за консервативним критерієм; на S2 додатково підтверджено стійкість до конкретної зміни зовнішньої сітки.

Водночас одночасна зміна жорсткості всіх трьох кілець не дозволяє встановити роль кожного кільця. Зокрема, від'ємний сумарний контраст не означає, що кожне з кілець окремо екранує тріщину. Це питання потребує розрахунків за вибіркової зміни матеріалу, з однаковою активною площею та незмінною геометрією. Йому присвячено наступний розділ. Питання зміни фактичної геометрії при незалежному чисельному розвитку тріщини розглянемо окремо.
